\section{Literature Review}
\label{sec:lit}

We review research incorporating reasoning and critique data into fine-tuning, building toward the critique-augmented preference optimisation approach of this work. We adopt a unified notation: $\pi_\theta$ denotes the learnable policy, $\pi_{\text{ref}}$ the frozen reference (usually the base SFT model), $x$ the prompt, $y_w$ and $y_l$ the preferred and rejected responses, $c$ a model's own chain-of-thought reasoning trace, $c_v$ an oracle or verifier chain-of-thought critique (``CoT of the verifier''), $r_\phi(x,y)$ a learned scalar reward, and $\hat{r}_\theta(x,y) = \beta \log \frac{\pi_\theta(y|x)}{\pi_{\text{ref}}(y|x)}$ the implicit reward.

The review is organised around a central question: \textit{what kind of signal is being used, and for which model?} The narrative moves from binary AI feedback (Section~\ref{sec:cot-motivation}), through enriching DPO with supervised reasoning targets (Section~\ref{sec:dpo-nll}), to methods that improve verifiers by having them reason (Section~\ref{sec:gen-verify}), and finally to the closest work --- methods that bring verifier critique directly into the base model's training objective (Section~\ref{sec:base-critique}).

% ============================================================
\subsection{Beyond Binary Feedback}
\label{sec:cot-motivation}
% ============================================================

\subsubsection{RLAIF and the Information Bottleneck}

RL from AI Feedback (RLAIF) replaces human annotators with LLM oracles. \autocite{bai2022constitutionalaiharmlessnessai} first use LLM-generated labels for harmlessness finetuning, doing RL with a preference model trained with LLM preferences. \autocite{lee2024rlaifvsrlhfscaling} use LLM to directly give rewards, showing that RLAIF matches or exceeds RLHF on summarisation and dialogue, even when the labeller is the same model being aligned. \autocite{li2024hrlaifimprovementshelpfulnessharmlessness} instruct the labeller to reason explicitly and report modest improvements. A shared limitation is that the oracle's reasoning is discarded after labelling: the sole training signal is a binary or scalar preference label, creating an information bottleneck that the remainder of this review seeks to address.

\subsubsection{Impact of Chain-of-Thought Data for DPO}
\label{sec:cot-dpo}

\autocite{gao2025uncoveringimpactchainofthought} provide a systematic empirical study of CoT's impact on DPO. In text-to-SQL --- where preference datasets traditionally contain only final SQL queries \textit{without} reasoning steps --- DPO degrades performance across the majority of ten tested LLMs. Augmenting the same datasets with synthetically generated chain-of-thought solutions achieves consistent DPO improvements for the first time. The analysis identifies three mechanisms. First, \textit{reward hacking mitigation}: without CoT, the implicit reward \eqref{eq:implicit_reward} can be inflated by superficial correlates (longer SQL, non-executable patterns) that do not reflect reasoning quality. Model performance initially peaks then decreases, while the implicit reward continues to rise. With CoT, the model's self-assessed implicit reward tracks actual performance. Second, \textit{discriminative calibration}: CoT-trained models maintain stable classification accuracy between correct and incorrect solutions throughout training, whereas vanilla DPO becomes unstable. Third, \textit{sample efficiency}: CoT-augmented DPO continues to improve with additional data, while vanilla DPO saturates. 2 samples of CoT with DPO achieves the same performance gains as 16 samples of vanilla DPO.

This finding directly motivates our approach: if the model's own CoT improves DPO, then an external verifier's critique --- which provides a richer perspective on why one response is preferred --- should be informative too.

\subsubsection{Self-Taught Reasoning: STaR and V-STaR}

\textbf{STaR.} \autocite{zelikman2022starbootstrappingreasoningreasoning} establish that CoT data, even when self-generated, is a richer training signal than outcome labels alone. Given a pretrained model $M$ and a dataset $\mathcal{D} = \{(x_i, y_i)\}$ of problems with ground-truth answers, STaR begins with a small few-shot prompt set $\mathcal{P}$ containing rationale examples. At each iteration, the model generates a reasoning trace $\hat{c}_i$ and answer $\hat{y}_i$ for every problem, using $\mathcal{P}$ in the context. Only correct generations where $\hat{y}_i = y_i$ are retained; the model is then fine-tuned on this filtered set with the standard NLL objective on rationale-answer pairs:
\begin{equation}
    \mathcal{L}_{\text{STaR}}(M) = -\mathbb{E}_{(x,y) \sim \mathcal{D}} \left[\,\mathbf{1}[\hat{y} = y]\cdot \sum_{t} \log M(\hat{c}_t, \hat{y}_t \mid x, \hat{c}_{<t})\right]
    \label{eq:star}
\end{equation}

STaR with rationalization matches GPT-3 (30$\times$ larger) on CommonsenseQA while using only 86.7\% of training data. The method demonstrates that self-generated rationales -- even produced with $|\mathcal{P}| = 10$ few-shot examples -- provide a dramatically richer training signal than training on final answers alone.

\vspace{1em}
\textbf{V-STaR.} \autocite{hosseini2024vstartrainingverifiersselftaught} identify the core inefficiency of STaR: the large number of \emph{incorrect} self-generated solutions is discarded, despite containing information about error patterns. V-STaR repurposes these failures to co-train a generator $G$ and a verifier $V$ in parallel across $T$ iterations.

They maintain two accumulating datasets: $\mathcal{D}_{\text{GEN}}$: correct solutions $(x, \hat{y})$, used for generator SFT, and $\mathcal{D}_{\text{VER}}$: \emph{all} generated solutions $(x, \hat{y}, z)$, used for verifier training. At each iteration $t$, $k=16$ solutions per problem are sampled from $G^t$ and automatically labeled for correctness --- by comparing the extracted final answer against the ground-truth answer for math, or by running test cases for code. Correct solutions are appended to $\mathcal{D}_{\text{GEN}}$; all solutions are appended to $\mathcal{D}_{\text{VER}}$. $G^t$ is fine-tuned on $\mathcal{D}_{\text{GEN}}^{<t}$ with the standard NLL loss Eq.~\eqref{eq:star}. From $\mathcal{D}_{\text{VER}}$, preference pairs $(x, y^+, y^-)$ are formed by taking the pair of all correct solutions $y^+$ and all incorrect solutions $y^-$ for each problem $x$. The verifier is trained with DPO using $G_{\text{SFT}}$ (the base model fine-tuned once) as the fixed reference policy:
\begin{equation}
    \mathcal{L}_{\text{DPO}}(V^t;\, G_{\text{SFT}}) = -\mathbb{E}_{(x, y^+, y^-) \sim \mathcal{D}_{\text{VER}}} \!\left[\log \sigma\!\left(\hat{r}(x, y^+) - \hat{r}(x, y^-)\right)\right], \quad \hat{r}(x, y) = \beta \log \frac{V^t(y \mid x)}{G_{\text{SFT}}(y \mid x)}
    \label{eq:vstar}
\end{equation}

At inference time, the verifier is used in Best-of-k reranking to choose solutions for better performance. Crucially, the verifier provides no training-time signal to the generator. The generator is trained purely via SFT on correct solutions (identical to STaR). The verifier's direct contribution is better performance at inference via Best-of-$k$ reranking. As the generator improves across iterations it produces harder-to-distinguish incorrect solutions, enriching $\mathcal{D}_{\text{VER}}$. However this does not feed any gradient back into the generator's training. This distinguishes V-STaR from later methods. Nevertheless, V-STaR 7B surpasses LLaMA-2 70B (8-shot) on GSM8K and nearly matches CodeLLaMA 34B on HumanEval in out-of-domain transfer, which illustrates the value of CoT training signals.

% ============================================================
\subsection{DPO with NLL}
\label{sec:dpo-nll}
% ============================================================

A recurring structural pattern in the literature combines the DPO preference loss with an auxiliary negative log-likelihood (NLL) term on a supervised target. We present the generic template first, then show how each paper instantiates it with a \textbf{different NLL target} $t$:
\begin{equation}
    \mathcal{L}_{\text{DPO+NLL}} = \underbrace{-\log \sigma\!\left(\hat{r}_\theta(x, y_w) - \hat{r}_\theta(x, y_l)\right)}_{\text{preference term}} \;+\; \alpha \underbrace{\mathcal{L}_{\text{NLL}}(t \mid x)}_{\text{supervised anchor}}
    \label{eq:dpo_nll_template}
\end{equation}
The DPO term provides a contrastive signal (push apart $y_w$ and $y_l$), while the NLL term anchors the policy to a specific supervised target. As we show below, the choice of $t$ is what distinguishes these methods from one another.

\subsubsection{Iterative Reasoning Preference Optimization (IRPO)}
\label{sec:irpo}

\autocite{pang2024iterativereasoningpreferenceoptimization} set the NLL target to the \textit{model's own winning chain-of-thought}: $t = (c_w, y_w)$. At each iteration, the model samples $N$ reasoning traces per prompt, partitions them into correct ($c_w, y_w$) and incorrect ($c_l, y_l$) by checking the final answer against ground truth, and optimises:
\begin{align}
    \mathcal{L}_{\text{IRPO}} & = \mathcal{L}_{\text{DPO}}(c_w, y_w, c_l, y_l \mid x) + \alpha\,\mathcal{L}_{\text{NLL}}(c_w, y_w \mid x) \nonumber                                                                                                                                              \\
                              & = -\log \sigma\!\left(\beta \log \frac{\pi_\theta(c_w, y_w \mid x)}{\pi_{\text{ref}}(c_w, y_w \mid x)} - \beta \log \frac{\pi_\theta(c_l, y_l \mid x)}{\pi_{\text{ref}}(c_l, y_l \mid x)}\right) - \alpha \frac{\log \pi_\theta(c_w, y_w \mid x)}{|c_w| + |y_w|}
    \label{eq:irpo}
\end{align}
The NLL term is length-normalised by $|c_w| + |y_w|$ to avoid a bias toward shorter sequences. Crucially, the reference model $\pi_{\text{ref}}$ is updated to $\pi_\theta$ at each iteration before re-sampling. Under standard DPO (i.e.\ $\alpha = 0$), the log-probabilities of \emph{chosen} sequences actually decrease during training --- the model achieves the required preference margin by pushing rejected probabilities down faster than chosen ones fall. The NLL anchor prevents this collapse: with $\alpha > 0$, chosen log-probabilities increase while rejected ones decrease, creating a healthy expanding margin. Empirically, DPO alone achieves 61.8\% on GSM8K while DPO+NLL achieves 73.1\% at iteration~1; after four iterations.

\subsubsection{AIPO: Agreement-Aware Iterative Preference Optimization}
\label{sec:aipo}

\autocite{shen2024aipoimprovingtrainingobjective} set the NLL target to the \textbf{winning response}: $t = y_w$. Their key modification is to the DPO term itself. Defining the log-probability ratios
$$
    s_\theta = \log \frac{\pi_\theta(y_w \mid x)}{\pi_\theta(y_l \mid x)}, \qquad s_{\text{ref}} = \log \frac{\pi_{\text{ref}}(y_w \mid x)}{\pi_{\text{ref}}(y_l \mid x)},
$$
standard DPO is $-\log \sigma(\beta(s_\theta - s_{\text{ref}}))$. AIPO introduces an agreement-aware coefficient $\alpha > 0$:
\begin{equation}
    \mathcal{L}_{\alpha\text{-DPO}} = -\mathbb{E}_{\mathcal{D}}\!\left[\log \sigma\!\left(\beta\big(s_\theta - (1 + \alpha)\,s_{\text{ref}}\big)\right)\right]
    \label{eq:alpha_dpo}
\end{equation}
The gradient weight decomposes as $\sigma(\beta(s_{\text{ref}} - s_\theta + \alpha\, s_{\text{ref}}))$: the additional $\alpha\,s_{\text{ref}}$ amplifies the learning signal when the reference model agrees with the preference label and dampens it when there is disagreement, mitigating the length exploitation that compounds across iterative DPO. The full AIPO loss combines this with a length-normalised NLL term:
\begin{equation}
    \mathcal{L}_{\text{AIPO}} = \mathcal{L}_{\alpha\text{-DPO}} + \lambda \cdot \mathcal{L}_{\text{NLL}}, \qquad \mathcal{L}_{\text{NLL}} = -\frac{1}{|y_w|}\log \pi_\theta(y_w \mid x)
    \label{eq:aipo}
\end{equation}
AIPO achieves 57.1\% length-controlled win rate on AlpacaEval~2.0 with Gemma-2-9B-It and 81.6\% on Arena-Hard with Mistral-Large, outperforming GPT-4.

\subsubsection{Direct Judgement Preference Optimization (DJPO)}
\label{sec:djpo}

\autocite{wang2024directjudgementpreferenceoptimization} set the NLL target to the \textbf{verifier's CoT critique and judgment}: $t = (c_v, j)$, where $c_v$ is the verifier's chain-of-thought critique and $j$ is the final judgment. Crucially, the model being trained is a \textit{judge}, not the base policy. They construct three types of DPO preference pairs --- CoT critique $\mathcal{D}_{\text{CoT}}$, standard judgment $\mathcal{D}_{\text{Std}}$, and response deduction $\mathcal{D}_{\text{Ded}}$ --- and train the collective preference dataset with:
\begin{equation}
    \mathcal{L}_{\text{DJPO}} = -\frac{\log \pi_\theta(y_w^{} \mid x)}{|y_w| + |x|} - \log \sigma\!\left(\beta \log \frac{\pi_\theta(y_w \mid x)}{\pi_{\text{ref}}(y_w \mid x)} - \beta \log \frac{\pi_\theta(y_l \mid x)}{\pi_{\text{ref}}(y_l \mid x)}\right)
    \label{eq:djpo}
\end{equation}
where $y_w$ and $y_l$ are correct and incorrect judge evaluations respectively. Using 680K preference pairs with a 70:15:15 ratio across the three types, they train SFR-Judge models at 8B, 12B, and 70B scales. SFR-Judge-8B outperforms GPT-4o on aggregate pairwise evaluation across 7 benchmarks; SFR-Judge-70B achieves the best overall score of 84.25\%.

% ============================================================
\subsection{Generative Verifiers and Critique-Augmented Reward Models}
\label{sec:gen-verify}
% ============================================================

The methods above train the policy model. A parallel line of work improves the verifiers and reward models that judge responses, by having them produce reasoning before scoring. These methods establish that critique data is a powerful training signal for judging.

\subsubsection{Generative Verifiers (GenRM)}
\label{sec:genrm}

\autocite{zhang2024generativeverifiersrewardmodeling} observe that discriminative reward models require a separate classification head to assign scalar rewards and cannot use the LLM's generative capabilities --- they cannot reason step-by-step before scoring, and they cannot leverage additional inference-time compute. Generative verifiers (GenRM) address this by recasting verification as next-token prediction, training the verifier with the standard SFT loss. Gemini 1.0 Pro generates solutions for GSM8K while a Gemma model is fine-tuned as the verifier.

In the simplest variant (\textbf{GenRM-Direct}), given instruction $I = \text{``Is the answer correct (Yes/No)?''}$, the verifier is trained on:
\begin{equation}
    \mathcal{D}_{\text{Direct}} = \{(x, y^+, I),\; \text{`Yes'}\} \;\cup\; \{(x, y^-, I),\; \text{`No'}\}
    \label{eq:genrm_direct_data}
\end{equation}
where $y^+$ and $y^-$ are correct and incorrect solutions labelled by an external oracle. At inference, the verifier score is the next-token probability of `Yes':
\begin{equation}
    r_{\text{Direct}}(x, y) = \pi_\theta(\text{`Yes'} \mid x, y, I)
    \label{eq:genrm_direct_score}
\end{equation}
which is used with Best-of-N sampling to choose the best generator solution at inference time.

\textbf{Unifying generation and verification.} Training on correct solution generation provides positive transfer to verification. The default training objective for the verifier mixes both tasks:
\begin{equation}
    \mathcal{L}_{\text{GenRM}}(\theta) = \mathcal{L}_{\text{SFT}}(\theta, \mathcal{D}_{\text{verify}}) + \lambda\,\mathcal{L}_{\text{SFT}}(\theta, \mathcal{D}_{\text{correct}})
    \label{eq:genrm}
\end{equation}
where $\mathcal{D}_{\text{verify}}$ is either $\mathcal{D}_{\text{Direct}}$ or $\mathcal{D}_{\text{CoT}}$, and $\mathcal{D}_{\text{correct}}$ is a set of correct solutions from the training set.

\textbf{GenRM-CoT.} Since verification requires nuanced reasoning, GenRM-CoT inserts a chain-of-thought rationale $c_v$ before the final Yes/No token. The dataset is:
\begin{equation*}
    \mathcal{D}_{\text{CoT}} = \bigl\{(x, y^+, I_v),\,(c_v, I, \text{`Yes'})\bigr\} \cup \bigl\{(x, y^-, I_v),\,(c_v, I, \text{`No'})\bigr\}
\end{equation*}
where $I_v = \text{``Let's verify step by step''}$. The rationales $c_v$ are generated offline by a capable LLM. Rationales are filtered by majority agreement across multiple generated samples. The verifier itself generates no $c_v$ during training; only at inference does it sample its own rationale and then score:
\begin{equation}
    r_{\text{CoT}}(x, y) = \pi_\theta(\text{`Yes'} \mid x, y, I_v, c_v, I), \qquad c_v \sim \pi_\theta(\cdot \mid x, y, I_v)
    \label{eq:genrm_cot}
\end{equation}
Multiple independent rationales can be sampled and aggregated via majority voting to use additional test-time compute:
\begin{equation}
    r_{\text{MV}@K}(x,y) = \frac{1}{K}\sum_{i=1}^{K} \pi_\theta\!\left(\text{`Yes'} \mid x, y, I_v, c_v^{(i)}, I\right), \quad c_v^{(i)} \overset{\text{i.i.d.}}{\sim} \pi_\theta(\cdot \mid x, y, I_v)
    \label{eq:genrm_mv}
\end{equation}

The whole procedure is a single-pass offline SFT; the trained verifier is used only at inference for Best-of-$N$ reranking. GenRM-CoT achieves 94.2\% Best-of-32 on GSM8K (vs.\ $\sim$90\% for a discriminative RM), and generalises to MATH500 trained only on GSM8K with 6.4$\times$ better sample efficiency than the discriminative baseline. GenRM also outperforms DPO-trained verifiers substantially, with the best DPO scoring strategy (raw log-likelihood, without subtracting the reference) still underperforming GenRM-Direct.

\subsubsection{GenRM-STaR-DPO: Training Judges with DPO on Rationales}
\label{sec:genrm-star}

\autocite{mahan2024generativerewardmodels} (Mahan et al., 2024) propose a unified RLHF--RLAIF framework built on a key theoretical observation: under the Bradley-Terry model, the standard RLHF reward-difference objective is equivalent to a preference log-odds objective,
\begin{equation}
    r_\phi(x, y) - r_\phi(x, y_{\text{ref}}) = \log \frac{p_{\text{BT}}(y \succ y_{\text{ref}} \mid x)}{1 - p_{\text{BT}}(y \succ y_{\text{ref}} \mid x)}
    \label{eq:genrm_logodds}
\end{equation}
This means one can replace the scalar reward parameterisation with a fully general preference model $p_\phi(y_w \succ y_l \mid x)$ extracted directly from an LLM's token probabilities, without requiring a separate classification head or a specific distributional assumption. This is the conceptual foundation for training a generative judge as a reward model.

\textbf{Three training variants.} Given a preference dataset $\mathcal{D} = \{(x, y_1, y_2, I)\}$ where $I \in \{A, B\}$ is the winning response indicator:

\textbf{GenRM (no CoT)} --- standard SFT on the indicator token:
\begin{equation}
    \mathcal{L}_{\text{GenRM}}(\pi_\phi) = \mathbb{E}_{(x,y_1,y_2,I)\sim\mathcal{D}}\!\left[-\log \pi_\phi(I \mid x, y_1, y_2)\right]
    \label{eq:genrm_nochain}
\end{equation}

\textbf{CoT-GenRM with Rationalization} --- a post-rationalization model $q$ (e.g.\ GPT-4 or a stronger LLM) generates a critique rationale $c_v$ conditioned on the correct answer, and the model is trained on both:
\begin{equation}
    \mathcal{L}_{\text{Rationalizer}}(\pi_\phi) = \mathbb{E}_{(x,y_1,y_2,c_v,I)\sim\mathcal{D}}\!\left[-\log \pi_\phi(I \mid x, y_1, y_2, c_v) - \log \pi_\phi(c_v \mid x, y_1, y_2)\right]
    \label{eq:genrm_rat}
\end{equation}

\textbf{CoT-GenRM-STaR-DPO} --- instead of discarding incorrect self-generated critique rationales (as STaR-SFT does), STaR-DPO constructs preference pairs $(c_{v,w}, I_w)$ vs $(c_{v,l}, I_l)$ from correct and incorrect reasoning traces and applies DPO over the full rationale-plus-judgment sequence:
\begin{equation}
    \mathcal{L}_{\text{STaR-DPO}}(\pi_\phi) = -\mathbb{E}_{\mathcal{D}}\!\left[\log\sigma\!\left(\beta\log\frac{\pi_\phi(I_w, c_{v,w} \mid x, y_1, y_2)}{\pi_{\text{ref}}(I_w, c_{v,w} \mid x, y_1, y_2)} - \beta\log\frac{\pi_\phi(I_l, c_{v,l} \mid x, y_1, y_2)}{\pi_{\text{ref}}(I_l, c_{v,l} \mid x, y_1, y_2)}\right)\right]
    \label{eq:genrm_dpo}
\end{equation}
The model learns not just which response is correct but which \emph{reasoning patterns} lead to correct judgments, and this contrastive signal is what drives OOD generalisation.

On-policy rationales are essential: using rationales from a strong off-policy model (GPT-4 or Llama-3.1-70B) as the training signal --- the STaR-Rationalizer approach --- matches STaR-DPO in-distribution but collapses OOD (RewardBench accuracy drops to 63--70\%, below even the zero-shot LLM-as-a-Judge). The authors attribute this to distribution mismatch: the model learns to imitate the rationalizer's style in-distribution but generates different reasoning at test time. STaR-DPO avoids this by training on the model's own self-generated rationales throughout. STaR-DPO matches the best explicit reward models in-distribution (73.9\% vs 73-74\%) while outperforming all of them OOD, particularly on Safety (91.0\% vs 81.8\% for PairRM). On reasoning-specific tasks (UltraInteract training, RewardBench Reasoning), STaR-DPO achieves 87.2\% vs 70.8\% for the best non-reasoning model.

\subsubsection{Critique-out-Loud Reward Models (CLoud)}
\label{sec:cloud}

\autocite{ankner2024critiqueoutloudrewardmodels} train reward models that generate a natural-language critique $c$ before predicting a scalar reward $r_{\theta_R}(x, y, c)$, via a three-stage pipeline. First, the base model is SFT'd on oracle critiques (from Llama-3.1-405B). Second, oracle critiques are \emph{replaced} with self-generated ones to form $\mathcal{D}_{\text{self}}$. Third, the model is trained on the joint loss:
\begin{equation}
    \mathcal{L}_{\text{CLoud}} = \underbrace{-\mathbb{E}\!\left[\log\sigma\!\left(r_{\theta_R}(x, y_w, c_w) - r_{\theta_R}(x, y_l, c_l)\right)\right]}_{\text{reward modelling}} + \lambda \underbrace{\mathcal{L}_{\text{SFT}}(\theta_B, \theta_{LM},\, \mathcal{D})}_{\text{critique generation}}
    \label{eq:cloud}
\end{equation}
\textbf{Key insight: on-policy vs.\ off-policy critiques.} Training the reward head on oracle (off-policy) critiques causes a 5.6 percentage point drop in RewardBench accuracy for the 8B model compared to self-generated (on-policy) critiques, due to distribution shift between training and inference. This observation has direct implications for Method~1: the base model's NLL target $c_v$ should ideally be drawn from a distribution the model can learn to approximate at inference.

CLoud-8B improves pairwise classification accuracy by 4.65 points over classic reward models and achieves Pareto-dominant Best-of-N win rates on ArenaHard.

\subsubsection{Critic-RM: Critiques as Latent Variables}
\label{sec:critic-rm}

\autocite{yu2025selfgeneratedcritiquesboostreward} provide an elegant variational formulation that treats critiques as latent variables. Denoting $z_w, z_l$ as critiques for chosen and rejected responses, the preference likelihood is marginalised:
\begin{equation}
    \log p(y_w \!\succ\! y_l \mid x) \;\geq\; \mathbb{E}_{q_\phi(z_w | y_w, x),\, q_\phi(z_l | y_l, x)}\!\left[\log \frac{p(y_w \!\succ\! y_l,\, z_w, z_l \mid x)}{q_\phi(z_w | y_w, x)\, q_\phi(z_l | y_l, x)}\right]
    \label{eq:critic_rm}
\end{equation}
This variational lower bound decomposes into two objectives trained jointly: (1) a \textbf{preference modeling loss} $\ell_r = -\log \sigma(r_\psi(x, [y_w; z_w]) - r_\psi(x, [y_l; z_l]))$, where $[y; z]$ concatenates response and critique; and (2) a \textbf{critique generation loss} $\ell_c = -\frac{1}{K}\sum_{z \in Z} \log q_\phi(z \mid y, x)$, which trains the critique model to approximate the oracle distribution. A dynamic weight schedule $\lambda(t)$ emphasises critique generation in early epochs and gradually shifts to reward prediction.

Critic-RM improves reward modeling accuracy by 3.7--7.3\% on RewardBench. The 70B variant outperforms both GPT-4 and Llama-3.1-405B on critique quality, and matches full-data performance with only 10\% of training data.

% ============================================================
\subsection{Training the Base Model with Critique}
\label{sec:base-critique}
% ============================================================

Recall the three defining features of Method~1: \textbf{(A)} a DPO loss on answer preferences $(y_w, y_l)$; \textbf{(B)} an NLL term on the oracle/verifier's CoT critique $c_v$; and \textbf{(C)} applied to the base policy model. No single prior paper combines all three; the methods below each match two of the three.

% ============================================================
\subsubsection{Critique Fine-Tuning}
% ============================================================

\autocite{wang2025critiquefinetuninglearningcritique} ask: what happens if we use \emph{only} the NLL term on critique data, without any preference optimisation? In Critique Fine-Tuning (CFT), the base model is trained on the teacher's critique $c_v$ of a given query-response pair by maximising:
\begin{equation}
    \mathcal{L}_{\text{CFT}}(\theta) = -\sum_{t=1}^{|c_v|} \log \pi_\theta(c_{v,t} \mid c_{v,<t},\, [x;\, y])
    \label{eq:cft}
\end{equation}
where $[x;\, y]$ concatenates the query and a (potentially flawed) response, and $c_v$ is a critique generated by a teacher model (GPT-4o). At inference, the CFT-trained model generates direct answers without producing critiques.

The methods in Section~\ref{sec:gen-verify} improve reward models and judges by having them generate critique. A closely related cluster instead trains the base \emph{policy} model directly using critique signals. This section covers the methods most directly related to Method~1.

\subsubsection{SuperCorrect: Cross-Model Collaborative DPO}

\autocite{yang2024supercorrect} introduce a two-stage framework. Stage~1 distils hierarchical thought templates from a large teacher model (o1-preview) via SFT. Stage~2 applies \textbf{cross-model collaborative DPO} on teacher-guided error-correction traces. Given an erroneous student solution with the first error at step $k$, the teacher generates a correction trace $(a_k^+, c_k^+)$ (error analysis and correction guidance) while the student independently generates $(a_k^-, c_k^-)$. The DPO loss is:
\begin{equation}
    \mathcal{L}_{\text{Cross-DPO}} = -\mathbb{E}\!\left[\log\sigma\!\left(\beta\log\frac{\pi_\theta(a_k^+, c_k^+ \mid x, s_{<k})}{\pi_{\text{ref}}(a_k^+, c_k^+ \mid x, s_{<k})} - \beta\log\frac{\pi_\theta(a_k^-, c_k^- \mid x, s_{<k})}{\pi_{\text{ref}}(a_k^-, c_k^- \mid x, s_{<k})}\right)\right]
    \label{eq:supercorrect}
\end{equation}
where $s_{<k}$ are the correct preceding steps. The teacher's correction trace plays the role of $c_v$ in Method~1. SuperCorrect-7B achieves 70.2\% on MATH (surpassing DeepSeekMath-7B by 7.8\%) and demonstrates that cross-model DPO at the step level substantially improves self-correction capabilities. CriticGPT \autocite{mcaleese2024llmcriticshelpllm} and CTRL \autocite{xie2025teachinglanguagemodelscritique} provide complementary evidence that LLMs can be explicitly taught to generate high-quality critiques, validating the practical feasibility of constructing the oracle critique data $c_v$ used in Method~1.

\subsubsection{Critique-GRPO: The Closest RL Analogue}
\label{sec:critique-grpo}

\autocite{zhang2025critiquegrpoadvancingllm} is the closest existing paper to Method~1 in spirit. Critique-GRPO integrates natural-language critiques with RL policy optimisation, training the base model to learn from both initial responses and critique-guided refinements. The combined objective is:
\begin{equation}
    J_{\text{C-GRPO}} = \frac{1}{n}\sum_{i=1}^{n}\sum_{t} r_t^{(i)}(\theta)\, \hat{A}_t^{(i)} \;+\; \frac{1}{k}\sum_{i'=1}^{k}\sum_{t} f\!\left(r_{\text{ref},t}^{(i')}(\theta)\right) \hat{A}_t^{(i')}
    \label{eq:critique_grpo}
\end{equation}
where $r_t^{(i)}(\theta) = \pi_\theta / \pi_{\theta_{\text{old}}}$ is the standard importance sampling ratio and $\hat{A}$ are group-normalised advantages computed over the mixed batch of initial and refined responses. The key technical contribution is a \textbf{policy shaping function} for refined responses:
$$
    f(x) = \frac{\pi_\theta}{\pi_\theta + \gamma}, \qquad 0 < \gamma < 1
$$
When $\pi_\theta$ is small (unfamiliar token from critique-guided refinement), $f \approx \pi_\theta / \gamma$ upweights the gradient; when $\pi_\theta$ is large (familiar token), $f \approx 1$. This amplifies learning from correct refinements that the model would not have explored autonomously. Removing the shaping function degrades average pass@1 from 47.08\% to 43.95\%.

\begin{center}
    \begin{tabular}{lll}
        \hline
                        & \textbf{Method~1}        & \textbf{Critique-GRPO} \\
        \hline
        Algorithm       & DPO (offline)            & GRPO (online RL)       \\
        Critique source & External oracle          & Self-generated         \\
        Training regime & Offline preference pairs & Online rollouts        \\
        \hline
    \end{tabular}
\end{center}

Critique-GRPO achieves $+$4.4\% on Qwen2.5-7B-Base and $+$16.7\% pass@1 on AIME~2024 (self-critiquing). Both methods combine critique reasoning with preference optimisation for the base model, arriving at the same idea from different training paradigms.

\subsubsection{iGRPO: Self-Feedback Without Explicit Critique}

\autocite{hatamizadeh2026igrpo} extend GRPO with a two-stage self-feedback loop that outperforms Critique-GRPO \emph{without} generating natural-language critiques. In Stage~1, $N$ drafts are sampled and scored; the best draft $\hat{d}$ is selected. In Stage~2, the prompt is augmented with $\hat{d}$ and $G$ refinements are sampled for GRPO optimisation, with the total budget held constant ($N + G = G_{\text{GRPO}}$). The expected reward of the selected draft satisfies:
\begin{equation}
    \mathbb{E}[R(\hat{d})] = 1 - (1 - V_\theta(q))^N
    \label{eq:igrpo}
\end{equation}
which is monotonically increasing in the policy's success probability $V_\theta(q)$, creating a bootstrapping dynamic: as the policy improves, its best draft improves, which provides better conditioning for further refinement. iGRPO also delays entropy collapse, maintaining broader exploration.

Results: Nemotron-H-8B-Base (GRPO 41.1\% $\to$ Critique-GRPO 43.4\% $\to$ iGRPO \textbf{45.0\%}); state-of-the-art for 7B models on AIME~2024 (\textbf{85.6\%}) and AIME~2025 (\textbf{79.6\%}).

\subsubsection{RL Tango: Joint Generator--Verifier Training}

\autocite{zha2025rltangoreinforcinggenerator} simultaneously train both a policy generator $\pi_g$ and a generative process-level verifier $\pi_v$ using interleaved RL. The generator receives a \textbf{combined advantage} blending outcome and step-level signals from the verifier:
\begin{equation}
    \hat{A}_{g,t}^{(i)} = (1 - \alpha)\,\hat{A}_{g,\text{out},t}^{(i)} + \alpha\,\hat{A}_{g,\text{step},t}^{(i)}
    \label{eq:rl_tango}
\end{equation}
where $\hat{A}_{g,\text{step},t}^{(i)} = \sum_{k: I(k) \geq t} \frac{y_{\text{step},k} / K - \mu}{\sigma}$ accumulates normalised verifier judgments $y_{\text{step},k} \in \{-1, 1\}$ from all subsequent steps. The blending weight $\alpha$ \textbf{decays exponentially} during training, emphasising step-level guidance early and shifting to outcome signals to mitigate reward hacking. The verifier is trained with a class-aware advantage to handle the imbalance between correct and incorrect generator solutions, and is rewarded solely on outcome correctness --- its step-level CoT reasoning \emph{emerges} from outcome-level rewards alone.

Tango matches vanilla GRPO's 200-step accuracy in approximately 60 steps (a 3.3$\times$ efficiency improvement) and achieves best-in-class performance among 7B models, with the verifier leading on ProcessBench. The result validates the core premise of Method~1: verifier reasoning is a useful learning signal for the generator, even when the two are trained jointly from outcome rewards alone.

% ============================================================
\subsection{Positioning of Method~1}
\label{sec:positioning}
% ============================================================

The closest prior work to Method~1 divides into three groups that each match two of its three defining features:

\begin{center}
    \small
    \begin{tabular}{lccc}
        \hline
        \textbf{Paper}                                                   & \textbf{A: DPO on answers} & \textbf{B: NLL on verifier CoT} & \textbf{C: Base model} \\
        \hline
        \textbf{Method~1 (this paper)}                                   & \checkmark                 & \checkmark                      & \checkmark             \\
        IRPO \autocite{pang2024iterativereasoningpreferenceoptimization} & \checkmark                 & NLL on own CoT                  & \checkmark             \\
        AIPO \autocite{shen2024aipoimprovingtrainingobjective}           & \checkmark                 & NLL on $y_w$                    & \checkmark             \\
        DJPO \autocite{wang2024directjudgementpreferenceoptimization}    & \checkmark                 & \checkmark                      & $\times$ (judge)       \\
        GenRM-STaR-DPO \autocite{mahan2024generativerewardmodels}        & \checkmark                 & \checkmark                      & $\times$ (judge)       \\
        CFT \autocite{wang2025critiquefinetuninglearningcritique}        & $\times$                   & \checkmark                      & \checkmark             \\
        Critique-GRPO \autocite{zhang2025critiquegrpoadvancingllm}       & GRPO                       & Self-critique                   & \checkmark             \\
        SuperCorrect \autocite{yang2024supercorrect}                     & On traces                  & Teacher correction              & \checkmark             \\
        RL Tango \autocite{zha2025rltangoreinforcinggenerator}           & Via RL                     & Verifier CoT                    & Co-training            \\
        \hline
    \end{tabular}
\end{center}

Method~1 sits at the intersection of three lines of work that have so far only been combined in pairs: (1) IRPO's DPO+NLL loss structure --- the right loss form, but the wrong NLL target; (2) DJPO and GenRM-STaR-DPO's use of verifier rationale as the NLL target --- the right data, but training the wrong model; and (3) CFT's base-model orientation --- the right model, but missing the DPO preference term. Concurrently, Critique-GRPO arrives at a closely analogous idea from the online RL direction, providing independent empirical validation that joint critique-reasoning and preference optimisation is a productive direction for base model training.
